% Options for packages loaded elsewhere
\PassOptionsToPackage{unicode}{hyperref}
\PassOptionsToPackage{hyphens}{url}
%
\documentclass[
  ignorenonframetext,
]{beamer}
\usepackage{pgfpages}
\setbeamertemplate{caption}[numbered]
\setbeamertemplate{caption label separator}{: }
\setbeamercolor{caption name}{fg=normal text.fg}
\beamertemplatenavigationsymbolsempty
% Prevent slide breaks in the middle of a paragraph
\widowpenalties 1 10000
\raggedbottom
\setbeamertemplate{part page}{
  \centering
  \begin{beamercolorbox}[sep=16pt,center]{part title}
    \usebeamerfont{part title}\insertpart\par
  \end{beamercolorbox}
}
\setbeamertemplate{section page}{
  \centering
  \begin{beamercolorbox}[sep=12pt,center]{part title}
    \usebeamerfont{section title}\insertsection\par
  \end{beamercolorbox}
}
\setbeamertemplate{subsection page}{
  \centering
  \begin{beamercolorbox}[sep=8pt,center]{part title}
    \usebeamerfont{subsection title}\insertsubsection\par
  \end{beamercolorbox}
}
\AtBeginPart{
  \frame{\partpage}
}
\AtBeginSection{
  \ifbibliography
  \else
    \frame{\sectionpage}
  \fi
}
\AtBeginSubsection{
  \frame{\subsectionpage}
}
\usepackage{amsmath,amssymb}
\usepackage{iftex}
\ifPDFTeX
  \usepackage[T1]{fontenc}
  \usepackage[utf8]{inputenc}
  \usepackage{textcomp} % provide euro and other symbols
\else % if luatex or xetex
  \usepackage{unicode-math} % this also loads fontspec
  \defaultfontfeatures{Scale=MatchLowercase}
  \defaultfontfeatures[\rmfamily]{Ligatures=TeX,Scale=1}
\fi
\usepackage{lmodern}
\ifPDFTeX\else
  % xetex/luatex font selection
\fi
% Use upquote if available, for straight quotes in verbatim environments
\IfFileExists{upquote.sty}{\usepackage{upquote}}{}
\IfFileExists{microtype.sty}{% use microtype if available
  \usepackage[]{microtype}
  \UseMicrotypeSet[protrusion]{basicmath} % disable protrusion for tt fonts
}{}
\makeatletter
\@ifundefined{KOMAClassName}{% if non-KOMA class
  \IfFileExists{parskip.sty}{%
    \usepackage{parskip}
  }{% else
    \setlength{\parindent}{0pt}
    \setlength{\parskip}{6pt plus 2pt minus 1pt}}
}{% if KOMA class
  \KOMAoptions{parskip=half}}
\makeatother
\usepackage{xcolor}
\newif\ifbibliography
\setlength{\emergencystretch}{3em} % prevent overfull lines
\providecommand{\tightlist}{%
  \setlength{\itemsep}{0pt}\setlength{\parskip}{0pt}}
\setcounter{secnumdepth}{-\maxdimen} % remove section numbering
\ifLuaTeX
  \usepackage{selnolig}  % disable illegal ligatures
\fi
\usepackage{bookmark}
\IfFileExists{xurl.sty}{\usepackage{xurl}}{} % add URL line breaks if available
\urlstyle{same}
\hypersetup{
  pdftitle={State of existing sampling in Mbalmayo and Bouamir observatories},
  pdfauthor={Gilles Dauby},
  hidelinks,
  pdfcreator={LaTeX via pandoc}}

\title{State of existing sampling in Mbalmayo and Bouamir observatories}
\author{Gilles Dauby}
\date{2024-06-03}

\begin{document}
\frame{\titlepage}

\begin{frame}{Mbalmayo / Bouamir / Loundoungou supersites /
observatories}
\phantomsection\label{mbalmayo-bouamir-loundoungou-supersites-observatories}
Two embedded permanent surveys

\begin{itemize}
\tightlist
\item
  Canopy trees censused by drone and field delineation
\item
  1 ha permanent plot where all trees \textgreater= 10 cm at breast
  height
\end{itemize}

\begin{longtable}{lrrrrr}
\toprule
 & \multicolumn{3}{c}{1 ha plots} & \multicolumn{2}{c}{Canopy observatory} \\ 
\cmidrule(lr){2-4} \cmidrule(lr){5-6}
site & Number of 
plots & Number of stem & Number of species & Number of trees & Number of species \\ 
\midrule\addlinespace[2.5pt]
Bouamir & 8 & 3689 & 218 & 1354 & 136 \\ 
Mbalmayo & 14 & 6865 & 243 & 2053 & 146 \\ 
\bottomrule
\end{longtable}
\end{frame}

\begin{frame}{(re)censuses}
\phantomsection\label{recensuses}
\begin{itemize}
\item
  In Mbalmayo, first three plots established in 06-2021. Recensus for 3
  plots in April 2024
\item
  Scheduled recensus in April 2025.
\item
  In Bouamir, first four plots established in 12-2018. Recensus in June
  2022.
\item
  Scheduled three to four additional plots in July 2024.
\end{itemize}
\end{frame}

\begin{frame}{Herbarium specimens}
\phantomsection\label{herbarium-specimens}
Taxonomic identification is difficult for most species. To confirm or
improve field identification, herbarium specimens should be collected
and identified in herbarium.

For both Bouamir and Mbalmayo, herbarium specimens are collected. 610
specimens are linked to individuals in plots, representing at least 220
different species.

We take on average 4 to 5 photos of each specimen, hence at least 2440
photos are available for feeding
\href{mailto:Pl@ntnet}{\nolinkurl{Pl@ntnet}}.
\end{frame}

\begin{frame}{Addtional information : weather stations}
\phantomsection\label{addtional-information-weather-stations}
\begin{itemize}
\item
  Mbalmayo : station at \textbf{ca.} 15 km to the observatory since
  September 2023
\item
  Bouamir : station within the observatory since 2017, thanks to
  collaboration with UCLA managing the station.
\end{itemize}
\end{frame}

\begin{frame}{Addtional information : weather stations}
\phantomsection\label{addtional-information-weather-stations-1}
Monthly mininmum and maximum temperature and rainfall and anomalies
since 2017 in Bouamir (Vincent Deblauwe, April 2024).

Small aside : strong anomalies in minimum temperature for the last 5
monts in Bouamir. El niño effect ?

\includegraphics[width=0.6\linewidth]{anomalies_climate_bouamir}
\end{frame}

\begin{frame}{Soil sampling}
\phantomsection\label{soil-sampling}
\begin{itemize}
\tightlist
\item
  Existing dataset for the four older plot in Bouamir (but only one
  sample per plot)
\end{itemize}

\begin{longtable}{lrrrrrr}
\toprule
plot\_name & ph\_soil & sand\_content & clay\_content & silt\_content & organic\_soil\_carbon & nitrogen\_soil\_carbon \\ 
\midrule\addlinespace[2.5pt]
bouamir001 & 3.99 & 37.90 & 45.96 & 16.14 & 2.02 & 0.19 \\ 
bouamir002 & 4.20 & 40.04 & 47.89 & 12.07 & 1.46 & 0.13 \\ 
bouamir003 & 4.53 & 38.04 & 47.89 & 14.07 & 1.07 & 0.10 \\ 
bouamir004 & 4.22 & 37.90 & 47.96 & 14.14 & 2.16 & 0.18 \\ 
\bottomrule
\end{longtable}
\end{frame}

\begin{frame}{Soil sampling}
\phantomsection\label{soil-sampling-1}
\begin{itemize}
\tightlist
\item
  Recently much more intense sampling took place for 6 plots in
  Mbalmayo. Sampling in Bouamir will take place in July.
\end{itemize}

\includegraphics[width=0.4\linewidth]{soil_sampling}

\begin{itemize}
\tightlist
\item
  Subgroup talk about soil can take place with David Bauman and myself
  and anybody interested (if anybody !)
\end{itemize}
\end{frame}

\begin{frame}{Soft traits sampling}
\phantomsection\label{soft-traits-sampling}
\begin{itemize}
\item
  For one year, seven field mission (six in Mbalmayo and one in
  Bouamir), we measured soft traits for canopy trees (as well as species
  collected in plots).
\item
  Climbers team (trained one year ago) to get branches from canopy
  (directly exposed to sunlight).
\item
  Objectives are to sample as many species as possible, starting from
  the most frequent.
\item
  Explore how traits covariation can be linked with phenological
  behaviours.
\end{itemize}
\end{frame}

\begin{frame}{Soft traits sampling}
\phantomsection\label{soft-traits-sampling-1}
In short, three leaves per cohort (growth units of different ages) are
sampled for LMA (at mesophyl, lamina and whole leaf scales), mesophyl
thickness, SPAD chlorophyll concentration, leaf area.

\includegraphics[width=0.1\linewidth]{cohort_trait_sampling}
\end{frame}

\begin{frame}{Soft traits sampling}
\phantomsection\label{soft-traits-sampling-2}
So far (ongoing work, some measurements still need to be processed) :

\begin{longtable}{lr}
\toprule
trait & n \\ 
\midrule\addlinespace[2.5pt]
SPAD\_chlorophyll\_measurement & 933 \\ 
SPAD\_chlorophyll\_measurement\_immature\_leaves & 82 \\ 
leaf\_area & 944 \\ 
lma\_lamina\_immature & 40 \\ 
lma\_lamina\_mature & 546 \\ 
lma\_whole\_leaf\_immature & 5 \\ 
lma\_whole\_leaf\_mature & 116 \\ 
mesophyl\_thickness & 922 \\ 
mesophyl\_thickness\_immature\_leaves & 81 \\ 
twig\_specific\_gravity & 260 \\ 
whole\_leaf\_area & 538 \\ 
\bottomrule
\end{longtable}
\end{frame}

\begin{frame}{Soft traits sampling}
\phantomsection\label{soft-traits-sampling-3}
\begin{itemize}
\item
  73 species have been sampled in \textbf{Mbalmayo}. 49.7\% of the total
  species pool, but species that represent 87\% of all canopy trees.
\item
  56 species have been sampled in \textbf{Bouamir}. 40.9\% of the total
  species pool, but species that represent 75\% of all canopy trees.
\end{itemize}
\end{frame}

\begin{frame}{Soft traits sampling}
\phantomsection\label{soft-traits-sampling-4}
\begin{itemize}
\tightlist
\item
  Recent internship by Vincent Pearlstein to study how traits
  (morphological, spectral and chemical properties) vary with leaf age,
  and ultimately better understand how photosynthetic capacity vary with
  leaf ages.
\end{itemize}

\includegraphics[width=0.4\linewidth]{vincent}
\includegraphics[width=0.4\linewidth]{vincent2}
\end{frame}

\begin{frame}{Soft traits sampling}
\phantomsection\label{soft-traits-sampling-5}
\begin{itemize}
\item
  5 species were selected in Mbalmayo. Several individuals (33) of each
  species were sampled several times during 3 months (February to
  April). Soft traits in addition to spectrometer measurements.
\item
  768 leaves measured, 13 114 spectras (4038 measurements).
\item
  Nitrogen and pigment concentrations will be available soon.
\item
  Estimate leaf age based on drone surveys, not easy (manual estimation
  based on RGB images).
\item
  Spectrometer used seemed to provide inconsistent measurements.
\end{itemize}
\end{frame}

\begin{frame}{Soft traits sampling}
\phantomsection\label{soft-traits-sampling-6}
Subgroup talk for more in depth explanation of the protocol and datasets
already available. Could help to select species for eco-physiological
measures.

We should make sure protocol are similar to have comparable datasets.

\includegraphics[width=0.4\linewidth]{leaf_schema}
\end{frame}

\begin{frame}{Data management (Work package 1)}
\phantomsection\label{data-management-work-package-1}
Gather, harmonize and share existing data to provide ground-truth
estimates of multiple dimensions of tree diversity in the Congo Bassin
Forests.

First of all, gather and organize data in a relational database to at
least make dataset comparable : the minimum is to have all data managed
with a similar taxonomic framework.
\end{frame}

\end{document}
